\documentclass[11pt]{article}
\usepackage[utf8]{inputenc}
\usepackage[T1]{fontenc}
\usepackage{geometry}
\geometry{a4paper, margin=1in}
\usepackage{enumitem}
\usepackage{titlesec}
\usepackage{fancyhdr}
\usepackage{lastpage}
\usepackage{tabularx}

% Header/Footer
\pagestyle{fancy}
\fancyhf{}
\rhead{Prompt Engineering - Makeup Exam}
\lhead{Ankara Medipol University}
\rfoot{Page \thepage\ of \pageref{LastPage}}

\title{\textbf{Prompt Engineering}\\\large Makeup Examination\\Fall 2025--2026}
\author{Instructor: Mehmet Ali Akyol, PhD\\Ankara Medipol University}
\date{}

\begin{document}

\maketitle

\noindent\textbf{Student Name:} \underline{\hspace{6cm}} \hfill \textbf{Student ID:} \underline{\hspace{4cm}}

\vspace{0.5cm}
\noindent\textbf{Instructions:}
\begin{itemize}[noitemsep]
    \item This exam consists of 20 multiple-choice questions.
    \item Each question is worth 5 points for a total of 100 points.
    \item Select the \textbf{best} answer for each question.
    \item You have 60 minutes to complete this exam.
    \item No electronic devices or notes are permitted.
    \item Mark your answers on the optical form provided.
\end{itemize}

\vspace{0.5cm}
\hrule
\vspace{0.5cm}

\begin{enumerate}

\item A \textbf{prompt} in the context of AI is best defined as:
\begin{enumerate}[label=\Alph*)]
    \item The internal memory of an LLM
    \item The training data used to build the model
    \item The hardware that runs the AI system
    \item The input message you give to an AI system to elicit a response
\end{enumerate}

\item What is \textbf{tokenization} in the context of LLMs?
\begin{enumerate}[label=\Alph*)]
    \item The process of encrypting user data
    \item Generating authentication tokens for API access
    \item Converting images into text descriptions
    \item Breaking text into smaller chunks (tokens) that the model can process
\end{enumerate}

\item What are \textbf{embeddings} in LLM architecture?
\begin{enumerate}[label=\Alph*)]
    \item Physical components of the computer running the model
    \item The user interface elements of AI applications
    \item Error logs generated during model training
    \item Numerical representations of tokens that capture semantic meaning
\end{enumerate}

\item The 2017 paper ``Attention Is All You Need'' introduced which key innovation?
\begin{enumerate}[label=\Alph*)]
    \item Self-attention mechanism (Transformer architecture)
    \item Tokenization
    \item Recurrent neural networks
    \item Convolutional neural networks
\end{enumerate}

\item \textbf{Hallucination} in LLMs refers to:
\begin{enumerate}[label=\Alph*)]
    \item The model refusing to answer questions
    \item The model asking clarifying questions
    \item The model generating plausible but incorrect or fabricated information
    \item The model running out of memory
\end{enumerate}

\item In the anatomy of a good prompt, which element specifies headings, JSON schema, or table columns?
\begin{enumerate}[label=\Alph*)]
    \item Task
    \item Context
    \item Constraints
    \item Output Specification
\end{enumerate}

\item \textbf{Zero-shot prompting} is best described as:
\begin{enumerate}[label=\Alph*)]
    \item Giving just the task without any examples
    \item Providing multiple examples before the task
    \item Using images along with text
    \item Chaining multiple prompts together
\end{enumerate}

\item \textbf{Few-shot prompting} is most useful when:
\begin{enumerate}[label=\Alph*)]
    \item You want to save tokens
    \item The task is extremely simple
    \item You need to steer style or format with specific examples
    \item You want to reduce response time
\end{enumerate}

\item Which is an example of a ``prompt anti-pattern''?
\begin{enumerate}[label=\Alph*)]
    \item Specifying output format clearly
    \item Providing relevant context
    \item Setting clear constraints
    \item Including conflicting directives like ``be brief'' and ``include all details''
\end{enumerate}

\item In Retrieval-Augmented Generation (RAG), what is the primary purpose?
\begin{enumerate}[label=\Alph*)]
    \item To speed up model inference
    \item To reduce the model size
    \item To ground answers with provided excerpts or references
    \item To train the model on new data
\end{enumerate}

\item The \textbf{CO-STAR} framework stands for:
\begin{enumerate}[label=\Alph*)]
    \item Code, Output, Style, Testing, Analysis, Review
    \item Context, Objective, Style, Tone, Audience, Response
    \item Constraint, Optimization, Structure, Target, Action, Result
    \item Chain, Order, Sequence, Task, Answer, Repeat
\end{enumerate}

\item \textbf{Meta-prompting} refers to:
\begin{enumerate}[label=\Alph*)]
    \item Using prompts to have an LLM generate or refine other prompts
    \item Prompting about metadata
    \item Using multiple models simultaneously
    \item Prompting in multiple languages
\end{enumerate}

\item Which dimension of linguistic register controls terminology density?
\begin{enumerate}[label=\Alph*)]
    \item Tenor (Relationship)
    \item Mode (Channel)
    \item Field (Topic Domain)
    \item Rhythm
\end{enumerate}

\item \textbf{Brand Voice Fingerprinting} involves:
\begin{enumerate}[label=\Alph*)]
    \item Creating biometric authentication
    \item Recording audio samples of brand representatives
    \item Extracting measurable voice attributes from existing content to ensure AI outputs match brand style
    \item Generating unique identifiers for users
\end{enumerate}

\item What is the main advantage of using ``LLM-as-a-Judge'' for evaluation?
\begin{enumerate}[label=\Alph*)]
    \item It is always 100\% accurate
    \item It eliminates the need for any human review
    \item It enables scalable evaluation of thousands of outputs quickly
    \item It reduces model training costs
\end{enumerate}

\item \textbf{Verbosity Bias} in LLM evaluation refers to:
\begin{enumerate}[label=\Alph*)]
    \item Models producing too few tokens
    \item Preference for technical jargon
    \item The tendency to rate longer answers as better, even if they contain fluff
    \item Bias toward English-language responses
\end{enumerate}

\item \textbf{Tree of Thoughts (ToT)} differs from Chain-of-Thought by:
\begin{enumerate}[label=\Alph*)]
    \item Using fewer tokens
    \item Being faster to execute
    \item Requiring no examples
    \item Exploring multiple reasoning branches and allowing backtracking
\end{enumerate}

\item The \textbf{IRAC} framework (Issue, Rule, Analysis, Conclusion) is specifically designed for:
\begin{enumerate}[label=\Alph*)]
    \item Medical diagnosis
    \item Financial reporting
    \item Software development
    \item Legal reasoning and analysis
\end{enumerate}

\item The term \textbf{Automation Bias} refers to:
\begin{enumerate}[label=\Alph*)]
    \item Bias in robotic systems
    \item Preference for automated testing
    \item The tendency for humans to over-rely on AI outputs without critical verification
    \item AI systems preferring automated tasks
\end{enumerate}

\item \textbf{Agentic AI} differs from traditional chatbots by being:
\begin{enumerate}[label=\Alph*)]
    \item Unable to use external tools
    \item Text-only without multimodal capabilities
    \item Proactive, multi-step, stateful, and capable of autonomous tool usage
    \item Slower and less accurate
\end{enumerate}

\end{enumerate}

\vspace{1cm}
\hrule
\vspace{0.5cm}

\begin{center}
\textbf{--- END OF EXAM ---}
\end{center}

\newpage
%=============================================================================
\section*{Optical Form Answer Sheet}
%=============================================================================

\noindent\textbf{Student Name:} \underline{\hspace{6cm}} \hfill \textbf{Student ID:} \underline{\hspace{4cm}}

\vspace{1cm}

\noindent\textbf{Instructions:} Fill in the circle completely for your answer choice. Use pencil only.

\vspace{1cm}

\begin{center}
\renewcommand{\arraystretch}{2}
\begin{tabular}{|c|c|c|c|c||c|c|c|c|c|}
\hline
\textbf{Q} & \textbf{A} & \textbf{B} & \textbf{C} & \textbf{D} & \textbf{Q} & \textbf{A} & \textbf{B} & \textbf{C} & \textbf{D} \\
\hline
1 & $\bigcirc$ & $\bigcirc$ & $\bigcirc$ & $\bigcirc$ & 11 & $\bigcirc$ & $\bigcirc$ & $\bigcirc$ & $\bigcirc$ \\
\hline
2 & $\bigcirc$ & $\bigcirc$ & $\bigcirc$ & $\bigcirc$ & 12 & $\bigcirc$ & $\bigcirc$ & $\bigcirc$ & $\bigcirc$ \\
\hline
3 & $\bigcirc$ & $\bigcirc$ & $\bigcirc$ & $\bigcirc$ & 13 & $\bigcirc$ & $\bigcirc$ & $\bigcirc$ & $\bigcirc$ \\
\hline
4 & $\bigcirc$ & $\bigcirc$ & $\bigcirc$ & $\bigcirc$ & 14 & $\bigcirc$ & $\bigcirc$ & $\bigcirc$ & $\bigcirc$ \\
\hline
5 & $\bigcirc$ & $\bigcirc$ & $\bigcirc$ & $\bigcirc$ & 15 & $\bigcirc$ & $\bigcirc$ & $\bigcirc$ & $\bigcirc$ \\
\hline
6 & $\bigcirc$ & $\bigcirc$ & $\bigcirc$ & $\bigcirc$ & 16 & $\bigcirc$ & $\bigcirc$ & $\bigcirc$ & $\bigcirc$ \\
\hline
7 & $\bigcirc$ & $\bigcirc$ & $\bigcirc$ & $\bigcirc$ & 17 & $\bigcirc$ & $\bigcirc$ & $\bigcirc$ & $\bigcirc$ \\
\hline
8 & $\bigcirc$ & $\bigcirc$ & $\bigcirc$ & $\bigcirc$ & 18 & $\bigcirc$ & $\bigcirc$ & $\bigcirc$ & $\bigcirc$ \\
\hline
9 & $\bigcirc$ & $\bigcirc$ & $\bigcirc$ & $\bigcirc$ & 19 & $\bigcirc$ & $\bigcirc$ & $\bigcirc$ & $\bigcirc$ \\
\hline
10 & $\bigcirc$ & $\bigcirc$ & $\bigcirc$ & $\bigcirc$ & 20 & $\bigcirc$ & $\bigcirc$ & $\bigcirc$ & $\bigcirc$ \\
\hline
\end{tabular}
\end{center}

\vspace{2cm}

\noindent\textbf{Instructor Signature:} \underline{\hspace{6cm}} \hfill \textbf{Date:} \underline{\hspace{4cm}}

\newpage
%=============================================================================
\section*{Answer Key (For Instructor Use Only)}
%=============================================================================

\begin{center}
\begin{tabular}{|c|c||c|c||c|c||c|c|}
\hline
\textbf{Q} & \textbf{Ans} & \textbf{Q} & \textbf{Ans} & \textbf{Q} & \textbf{Ans} & \textbf{Q} & \textbf{Ans} \\
\hline
1 & D & 6 & D & 11 & B & 16 & C \\
2 & D & 7 & A & 12 & A & 17 & D \\
3 & D & 8 & C & 13 & C & 18 & D \\
4 & A & 9 & D & 14 & C & 19 & C \\
5 & C & 10 & C & 15 & C & 20 & C \\
\hline
\end{tabular}
\end{center}

\vspace{1cm}

\section*{Detailed Answer Key with Explanations}

\begin{enumerate}
    \item \textbf{D} - A prompt is the input message given to an AI system.
    \item \textbf{D} - Tokenization breaks text into processable chunks.
    \item \textbf{D} - Embeddings are numerical vectors capturing semantic meaning.
    \item \textbf{A} - The Transformer architecture introduced self-attention.
    \item \textbf{C} - Hallucination is generating plausible but incorrect information.
    \item \textbf{D} - Output Specification defines format like JSON, tables.
    \item \textbf{A} - Zero-shot means no examples, just the task.
    \item \textbf{C} - Few-shot steers style/format with examples.
    \item \textbf{D} - Conflicting directives are a common anti-pattern.
    \item \textbf{C} - RAG grounds answers with provided context.
    \item \textbf{B} - CO-STAR: Context, Objective, Style, Tone, Audience, Response.
    \item \textbf{A} - Meta-prompting uses prompts to generate/refine prompts.
    \item \textbf{C} - Field (Topic Domain) controls terminology density.
    \item \textbf{C} - Brand Voice Fingerprinting extracts measurable attributes.
    \item \textbf{C} - LLM-as-Judge enables scalable evaluation.
    \item \textbf{C} - Verbosity Bias favors longer answers.
    \item \textbf{D} - ToT explores multiple branches with backtracking.
    \item \textbf{D} - IRAC is for Legal reasoning.
    \item \textbf{C} - Automation Bias is over-reliance on AI outputs.
    \item \textbf{C} - Agentic AI is proactive, multi-step, stateful.
\end{enumerate}

\vspace{1cm}

\section*{Topic Coverage by Question}

\begin{itemize}
    \item \textbf{Foundations} (Q1--Q5): Basic concepts of prompts, tokenization, embeddings, transformers, and hallucination
    \item \textbf{Prompt Design Principles} (Q6--Q10): Output specification, zero-shot, few-shot, anti-patterns, and RAG
    \item \textbf{Patterns and Frameworks} (Q11--Q12): CO-STAR framework and meta-prompting
    \item \textbf{Style and Format Control} (Q13--Q14): Linguistic register and brand voice fingerprinting
    \item \textbf{Evaluation Methods} (Q15--Q16): LLM-as-Judge and verbosity bias
    \item \textbf{Complex Problem Solving} (Q17): Tree of Thoughts
    \item \textbf{Domain-Specific Applications} (Q18): IRAC framework for legal reasoning
    \item \textbf{Ethics and Responsible AI} (Q19--Q20): Automation bias and agentic AI
\end{itemize}

\vspace{1cm}

\section*{Grading Scale}

\begin{center}
\begin{tabular}{|c|c|}
\hline
\textbf{Points} & \textbf{Grade} \\
\hline
90--100 & AA \\
85--89 & BA \\
80--84 & BB \\
75--79 & CB \\
70--74 & CC \\
65--69 & DC \\
60--64 & DD \\
50--59 & FD \\
0--49 & FF \\
\hline
\end{tabular}
\end{center}

\end{document}
